%Version 3.1 December 2024
% See section 11 of the User Manual for version history
%
%%%%%%%%%%%%%%%%%%%%%%%%%%%%%%%%%%%%%%%%%%%%%%%%%%%%%%%%%%%%%%%%%%%%%%
%%                                                                 %%
%% Please do not use \input{...} to include other tex files.       %%
%% Submit your LaTeX manuscript as one .tex document.              %%
%%                                                                 %%
%% All additional figures and files should be attached             %%
%% separately and not embedded in the \TeX\ document itself.       %%
%%                                                                 %%
%%%%%%%%%%%%%%%%%%%%%%%%%%%%%%%%%%%%%%%%%%%%%%%%%%%%%%%%%%%%%%%%%%%%%

%%\documentclass[referee,sn-basic]{sn-jnl}% referee option is meant for double line spacing

%%=======================================================%%
%% to print line numbers in the margin use lineno option %%
%%=======================================================%%

%%\documentclass[lineno,pdflatex,sn-basic]{sn-jnl}% Basic Springer Nature Reference Style/Chemistry Reference Style

%%=========================================================================================%%
%% the documentclass is set to pdflatex as default. You can delete it if not appropriate.  %%
%%=========================================================================================%%

%%\documentclass[sn-basic]{sn-jnl}% Basic Springer Nature Reference Style/Chemistry Reference Style

%%Note: the following reference styles support Namedate and Numbered referencing. By default the style follows the most common style. To switch between the options you can add or remove �Numbered� in the optional parenthesis. 
%%The option is available for: sn-basic.bst, sn-chicago.bst%  
 
%%\documentclass[pdflatex,sn-nature]{sn-jnl}% Style for submissions to Nature Portfolio journals
%%\documentclass[pdflatex,sn-basic]{sn-jnl}% Basic Springer Nature Reference Style/Chemistry Reference Style
\documentclass[pdflatex,sn-mathphys-num]{sn-jnl}% Math and Physical Sciences Numbered Reference Style
%%\documentclass[pdflatex,sn-mathphys-ay]{sn-jnl}% Math and Physical Sciences Author Year Reference Style
%%\documentclass[pdflatex,sn-aps]{sn-jnl}% American Physical Society (APS) Reference Style
%%\documentclass[pdflatex,sn-vancouver-num]{sn-jnl}% Vancouver Numbered Reference Style
%%\documentclass[pdflatex,sn-vancouver-ay]{sn-jnl}% Vancouver Author Year Reference Style
%%\documentclass[pdflatex,sn-apa]{sn-jnl}% APA Reference Style
%%\documentclass[pdflatex,sn-chicago]{sn-jnl}% Chicago-based Humanities Reference Style

%%%% Standard Packages
%%<additional latex packages if required can be included here>

\usepackage{graphicx}%
\usepackage{multirow}%
\usepackage{amsmath,amssymb,amsfonts}%
\usepackage{amsthm}%
\usepackage{mathrsfs}%
\usepackage[title]{appendix}%
\usepackage{xcolor}%
\usepackage{textcomp}%
\usepackage{manyfoot}%
\usepackage{booktabs}%
\usepackage{algorithm}%
\usepackage{algorithmicx}%
\usepackage{algpseudocode}%
\usepackage{listings}%
%%%%

%%%%%=============================================================================%%%%
%%%%  Remarks: This template is provided to aid authors with the preparation
%%%%  of original research articles intended for submission to journals published 
%%%%  by Springer Nature. The guidance has been prepared in partnership with 
%%%%  production teams to conform to Springer Nature technical requirements. 
%%%%  Editorial and presentation requirements differ among journal portfolios and 
%%%%  research disciplines. You may find sections in this template are irrelevant 
%%%%  to your work and are empowered to omit any such section if allowed by the 
%%%%  journal you intend to submit to. The submission guidelines and policies 
%%%%  of the journal take precedence. A detailed User Manual is available in the 
%%%%  template package for technical guidance.
%%%%%=============================================================================%%%%

%% as per the requirement new theorem styles can be included as shown below
\theoremstyle{thmstyleone}%
\newtheorem{theorem}{Theorem}%  meant for continuous numbers
%%\newtheorem{theorem}{Theorem}[section]% meant for sectionwise numbers
%% optional argument [theorem] produces theorem numbering sequence instead of independent numbers for Proposition
\newtheorem{proposition}[theorem]{Proposition}% 
%%\newtheorem{proposition}{Proposition}% to get separate numbers for theorem and proposition etc.

\theoremstyle{thmstyletwo}%
\newtheorem{example}{Example}%
\newtheorem{remark}{Remark}%

\theoremstyle{thmstylethree}%
\newtheorem{definition}{Definition}%

\raggedbottom
%%\unnumbered% uncomment this for unnumbered level heads




\def\blue#1{\textcolor[rgb]{0,0,1.0}{#1}}
\def\green#1{\textcolor[rgb]{0,0.6,0}{#1}}
\def\red#1{\textcolor[rgb]{1.0,0,0}{#1}}
\def\green#1{\textcolor[rgb]{0,0.6,0}{#1}}
\def\dgreen#1{\textcolor[rgb]{0.02,0.55,0.04}{#1}}
\def\gray#1{\textcolor[rgb]{0.5,0.5,0.7}{#1}}
\def\orange#1{#1}% reverted to black (was orange)
\def\Orange{}% reverted to black (was orange)

\def\Dgreen{\color[rgb]{0.02,0.55,0.04}}
\def\Blue{\color[rgb]{0,0,1}}
\def\Red{\color[rgb]{0,0,0}}% red removed (was red)
\def\Black{\color[rgb]{0,0,0}}
\def\Gray{\color[rgb]{0.5,0.5,0.7}}
\def\colorreset{\color[rgb]{0,0,0}}


\usepackage{xspace}

\def\ignore#1{\iffalse #1 \fi\xspace}
%\def\red#1{\RedColor\textcolor{red}{#1}}
\def\red#1{#1}% red removed (was \textcolor{red})
\def\blue#1{\textcolor{blue}{ #1}}
\def\ignore#1{\xspace}

\usepackage[normalem]{ulem}

\newcommand{\add}[1]{\dgreen{#1}}
%\newcommand{\erase}[1]{\textcolor{red}{\sout{\textcolor{black}{#1}}}}
\newcommand{\erase}[1]{\sout{#1}}% red removed (kept strikethrough)
% \rev: mark <k>=4 degree-aligned rerun revisions in red (use \long form so it spans \par)
\long\def\rev#1{{#1}}% red removed (was \color{red})

\def\subst#1#2{\add{#1}\erase{ (#2)}}

\def\replace#1#2{\dgreen{#2}}  % #1 old, #2 new
% \def\replace#1#2{\erase{#1} \dgreen{#2}}  % #1 old, #2 new

\renewcommand{\topfraction}{.99}
\renewcommand{\dbltopfraction}{.99}
\renewcommand{\bottomfraction}{.99}
\renewcommand{\textfraction}{.01}
\renewcommand{\floatpagefraction}{0.9}
\renewcommand{\dblfloatpagefraction}{0.9}


% parameters definition

\def\AgentSet{\mathcal{A}}
\def\post{\rho} % was P_p
\def\BoundedConfidence{\epsilon}
\def\e{\varepsilon}
\def\feed{F}

\def\maxfollower{S_{\it fol}}
\def\maxfeed{S_{\it fd}}
\def\Stubbornness{\mu_{\it stub}}
\def\ThofOpPreval{\varepsilon_s}


\usepackage{url}
\usepackage{subcaption}

\begin{document}

\title[A Mechanism of Structural Backfire Effect]{\dgreen{Mechanism of the Structural Backfire Effect: \\How Hub Opinion Shifts Fragment Heavy-Tailed Social Networks}}

%%=============================================================%%
%% GivenName	-> \fnm{Joergen W.}
%% Particle	-> \spfx{van der} -> surname prefix
%% FamilyName	-> \sur{Ploeg}
%% Suffix	-> \sfx{IV}
%% \author*[1,2]{\fnm{Joergen W.} \spfx{van der} \sur{Ploeg} 
%%  \sfx{IV}}\email{iauthor@gmail.com}
%%=============================================================%%

\author*[1]{\fnm{Tomoya} \sur{Takeda}}\email{tomoya.leo9@fuji.waseda.jp}

\author[2]{\fnm{Fujio} \sur{Toriumi}}\email{tori@sys.t.u-tokyo.ac.jp}

\author[1]{\fnm{Toshiharu} \sur{Sugawara}}\email{sugawara@waseda.jp}


\affil*[1]{\orgdiv{Department of Computer Science and Communications Engineering}, \orgname{Waseda University}, \orgaddress{\street{3-4-1 Okubo}, \city{Shinjyuku-ku}, \state{Tokyo}, \country{Japan}}}

\affil[2]{\orgdiv{Department of Systems Innovation}, \orgname{The University of Tokyo}, \\ \orgaddress{\street{7-3-1 Hongo}, \city{Bunkyo-ku}, \state{Tokyo}, \country{Japan}}}

%\affil[3]{\orgdiv{Department of Computer Science and Communications Engineering, Waseda University}, \orgname{Organization}, \orgaddress{\street{3-4-1 Okubo}, \city{Shinjyuku-ku}, \postcode{169-8555}, \state{Tokyo}, \country{Japan}}}

%%==================================%%
%% Sample for unstructured abstract %%
%%==================================%%

\abstract{
\Dgreen
Online discourse on social media has become a prerequisite for public
opinion formation for a large part of the population, and
influencers (or hub users) exert a strong influence on online communication.
We study a counterintuitive {\em structural backfire effect} in agent-based
opinion dynamics: when a moderate opinionated hub is driven toward an
extreme stance, it often fails to steer public opinion and instead accelerates
collective polarization by energizing opposite opinion followers.
We show that this is a robust {\em structural} phenomenon---emerging
specifically in heavy-tailed networks when users are sufficiently stubborn and
the opinion shifting hub is sufficiently central, and absent in networks lacking
a heavy tail. By directly measuring the causal chain, we identify its mechanism:
hub manipulation severs the primary cross partisan bridge, driving
opposite opinion followers to defect and collapse into echo chambers. 
The pre-manipulation model reproduces the empirical
stylized facts of real social media data, and the effect persists under a realistic gradual
opinion shift.
These findings advance our understanding of 
the vulnerability of social networks by highlighting 
how nonlinear structural dynamics produce counterintuitive collective
outcomes, 
thereby contributing to 
the establishment of conditions for sound web communication.
\Black
}
\keywords{Opinion Dynamics, Computational Social Science, Multi-Agent Simulation, Polarization, Influence Operation, Social Media}

%%================================%%
%% Sample for structured abstract %%
%%================================%%

% \abstract{\textbf{Purpose:} The abstract serves both as a general introduction to the topic and as a brief, non-technical summary of the main results and their implications. The abstract must not include subheadings (unless expressly permitted in the journal's Instructions to Authors), equations or citations. As a guide the abstract should not exceed 200 words. Most journals do not set a hard limit however authors are advised to check the author instructions for the journal they are submitting to.
% 
% \textbf{Methods:} The abstract serves both as a general introduction to the topic and as a brief, non-technical summary of the main results and their implications. The abstract must not include subheadings (unless expressly permitted in the journal's Instructions to Authors), equations or citations. As a guide the abstract should not exceed 200 words. Most journals do not set a hard limit however authors are advised to check the author instructions for the journal they are submitting to.
% 
% \textbf{Results:} The abstract serves both as a general introduction to the topic and as a brief, non-technical summary of the main results and their implications. The abstract must not include subheadings (unless expressly permitted in the journal's Instructions to Authors), equations or citations. As a guide the abstract should not exceed 200 words. Most journals do not set a hard limit however authors are advised to check the author instructions for the journal they are submitting to.
% 
% \textbf{Conclusion:} The abstract serves both as a general introduction to the topic and as a brief, non-technical summary of the main results and their implications. The abstract must not include subheadings (unless expressly permitted in the journal's Instructions to Authors), equations or citations. As a guide the abstract should not exceed 200 words. Most journals do not set a hard limit however authors are advised to check the author instructions for the journal they are submitting to.}

%%\pacs[JEL Classification]{D8, H51}

%%\pacs[MSC Classification]{35A01, 65L10, 65L12, 65L20, 65L70}

\maketitle

%\Blue

\section{Introduction}\label{Intro}
The advent of social media platforms such as Facebook, X (formerly Twitter), YouTube, and TikTok has profoundly transformed the global information environment. By significantly reducing communication barriers, these platforms offer novel interaction modes in which finding like-minded communities is effortless and the cost of broadcasting to the public has drastically decreased~\cite{Cheung2011,Carr2015,Nyhan2023}. Owing to their accessibility and efficiency, online social networks are increasingly envisioned as essential infrastructure for civic deliberation and digital democracy~\cite{Smith2021,Tseng2022,Yang2025}. However, to fully harness the potential of these transformative technologies, we must understand the complex interplay between system features and user behavior, particularly because the relatively short history of human interaction with social media means that our understanding of its long-term societal impact remains nascent.
\par


Despite these advantages, society faces significant challenges, including opinion polarization, misinformation diffusion, and malicious influence operations~\cite{Fisher2015,Baccarella2018,Conover2011}. Although phenomena such as polarization are not exclusive to the digital age, the architectural and functional characteristics of social media arguably catalyze their emergence. Structural phenomena, such as \emph{echo chambers}~\cite{Cinelli2020,Sasahara2021}, in which users are isolated in homophilic clusters driven by their own selective tie formation, and \emph{filter bubbles}~\cite{Uthsav2020,Lien2022}, driven by algorithmic filtering, create environments that foster radicalization. 
\Dgreen
A less studied yet equally consequential source of polarization 
is the opinion shift of highly connected \emph{hub users}, 
whose stance changes can propagate rapidly through the network. 
Such shifts arise spontaneously in events such as the sudden political revelations of
prominent influencers, the ideological radicalization of public figures,
or the adoption of increasingly extreme positions by populist politicians.
Such events routinely trigger follower defection and affective polarization~\cite{Flamino2023,Bail2018}.
Exogenous \emph{influence operations}~\cite{Fisher2015,Pascal2016,Golovchenko2020},
which deliberately target opinion leaders to steer collective discourse,
represent an intentional variant of the same mechanism.
\par


Whether endogenous or exogenous, the structural and behavioral conditions under which
such hub-driven opinion shifts trigger irreversible polarization remain 
imperfectly understood. 
While prior research has repeatedly shown that attempts to shift opinions 
through social influence frequently produce counterintuitive and nonlinear 
outcomes, including the reinforcement of prior positions rather than convergence
toward the intended target~\cite{Brehm1966,Bail2018}, the network-structural 
conditions modulate these hub-mediated dynamics has not been fully explained. 
In preliminary work, we proposed an {\em agent-based model} (ABM), which integrates well-established
social interaction models with behavioral repulsion (tie severing through
unfollowing rather than attitudinal reactance), and reported a novel
network-level phenomenon termed the {\em structural backfire effect}, in which the opinion
shift of moderately opinionated hub users toward extremes severs the cross-partisan bridge they
provide, driving far-opinion followers to defect and retreat into echo chambers,
ultimately intensifying polarization~\cite{Takeda2026}.
We use \emph{structural} to distinguish this network-mediated effect from the
classical, individual-level ``backfire effect,'' in which exposure to opposing
views reinforces a person's prior \emph{opinion}~\cite{Brehm1966,Nyhan2021}:
here, the intervention instead amplifies the opposing camp's \emph{posting activity}
through structural fragmentation of the network, with no opinion-level repulsion
term in the model. Hereafter we refer to it as the structural backfire effect, or
simply backfire where unambiguous.
However, these preliminary findings were limited to a specific parameter set,
a single network topology, and two hub nodes undergoing instantaneous opinion shifts;
moreover, the model lacked empirical grounding to establish its behavioral credibility
relative to real social media dynamics.
It therefore remained unclear whether the backfire effect is a universal structural
phenomenon or an artifact arising only under narrow conditions.
\par


This paper establishes the structural backfire effect as a robust 
structural phenomenon and explains its mechanism through four main contributions.
First, we provide direct mechanistic evidence for the causal chain underlying the backfire
effect by tracking unfollowing events, follower composition, and feed homogeneity at each
manipulation step. This reveals echo chamber deepening via 
hub loss as a structural mechanism for backfire, which is 
distinct from the psychological reactance processes identified in prior literature~\cite{Bail2018}.
Second, we identify the boundary conditions under which hub-driven opinion shifts trigger irreversible
polarization through a fine-grained parameter sweep across agent stubbornness and
opinion prevalence thresholds evaluated in heavy-tailed, high-clustering, and random networks.
Third, we validate the model against empirical stylized facts drawn from Twitter/X data,
establishing its behavioral credibility relative to real social media dynamics prior to the main analysis.
Fourth, in contrast to the instantaneous opinion shift modeled in our preliminary study,
we adopt a gradual shift model in which the hub's opinion is incrementally displaced
per platform access, demonstrating that the backfire effect persists under this
more realistic behavioral assumption.
Our results establish the backfire effect as a robust structural phenomenon consistently
associated with heavy-tailed degree distributions and high user stubbornness,
providing a theoretical framework that highlights the
critical role of hub users in network resilience and echo chamber formation.
\Black


\section{Related Work}\label{RelatedWork}
Computational social system modeling using ABM has long been a
primary method for gaining insight 
in complex social phenomena. Early research assumed individuals
as nodes in a graph with internal state variables and analyzed system
convergence and dynamics through local interactions~\cite{Clifford1973,Liggett1975,Axelrod1997,Sznajd2005,Galam2008}. 
Liggett et al.~\cite{Liggett1975} proposed the \emph{voter model}, in which agents
hold binary values and probabilistically change their internal values
adopting the state of a neighbor. This model serves as a foundational
framework to understand public opinion dynamics as a diffusion
processes on social networks.
\par


These discrete models were extended by incorporating greater
sociological realism~\cite{DeGroot1974,Friedkin1990,Deffuant2000,Hegselmann2002}. DeGroot~\cite{DeGroot1974} introduced a weighted
averaging model, in which an agent's opinion is updated based on the
weighted average of their neighbors' opinions, representing the
social influence process. Building on this concept, Deffuant et al.~\cite{Deffuant2000} and Hegselmann and Krause~\cite{Hegselmann2002}
introduced the concept of a \emph{bounded confidence model}, which
implies that agents interact only with those whose opinions fall
within a specified threshold range. This mechanism accounts for 
the psychological reality that individuals are unlikely to be influenced
by divergent views, thereby providing a mathematical basis for the emergence
of a local consensus or cluster.
\par



Social media has introduced novel forms of communication
environments that are distinct from traditional face-to-face interactions,
characterized by features such as algorithmic recommender systems and
dynamic tie formation (i.e., following and unfollowing
relationships). The complex dynamics that emerge from these features pose
critical questions that must be fully elucidated; therefore, a
vast body of literature has emerged to investigate these
phenomena~\cite{Watts2007,Lazer2009,Golder2011,Bond2012,Bakshy2015,DelVicario2016,Vosoughi2018,Sasahara2021,Conover2011}. Vosoughi
et al.~\cite{Vosoughi2018} conducted a large-scale analysis on
information diffusion on Twitter, and found that falsehoods spread more
broadly, rapidly, and deeply than truths. They attribute this pattern 
primarily to the greater novelty of false information, which increases
users’ propensity to share such information. Furthermore, evidence indicates
that humans, not automated accounts, serve as principal drivers
of falsehood diffusion. Sasahara et al.~\cite{Sasahara2021}
demonstrated that the interplay between social influence mechanisms
and the unfollowing of social ties accelerates polarization, thereby facilitating
the formation of echo chambers. The authors validated their simulations
against empirical data from X (formerly Twitter) regarding
political polarization during elections~\cite{Conover2011}, thus confirming
qualitative alignment between the model and real-world data. Therefore,
the information era has enabled researchers to leverage large-scale
behavioral logs for comprehensive analysis and model validation. This integration of
big-data analytics and multi-agent simulation grounded in
social science theories crystallized into the interdisciplinary
field of \emph{computational social science} (CSS)~\cite{Lazer2009}.
\par


In social network analysis, extensive research has
focused on the role of influencers and hub nodes in social networks. A
representative domain is \emph{influence maximization}, which
addresses the problem of selecting seed nodes to maximize the spread
of information (i.e., ``word-of-mouth'') under resource constraints,
primarily in the viral marketing context~\cite{Kempe2003,Kitsak2010,Budak2011,Kumar2022}. 
Kempe et al.~\cite{Kempe2003} provided the first efficient algorithms with
provable approximation guarantees for the influence maximization
problem, which is defined as the NP-hard task of identifying a set of nodes to
trigger the largest adoption cascade. More recently, Kumar et
al.~\cite{Kumar2022} employed graph embeddings and {\em graph neural networks} to efficiently solve this problem, demonstrating that deep-learning approaches outperformed classical heuristics.
\par


\Dgreen
Although these studies treated hub nodes as static conduits for information spread,
they did not address the consequences of hub users' \emph{own} opinion shifts, whether
endogenous, as discussed in Section~\ref{Intro}, or exogenously orchestrated through
\emph{influence operations}~\cite{Beskow2019,Ross2019,Bergh2020,Goldstein2023} such
as bot-driven campaigns, strategic promotion, account hijacking, and cognitive warfare.
Even within the influence operations literature, the emphasis has been on the
attacking actors rather than on the network position of the targeted hubs:
Ross et al.~\cite{Ross2019}, for instance, found that the number and sophistication
of bots outweigh their network position as drivers of opinion influence.
As a result, few studies have systematically examined the network-structural
conditions under which the opinion shift of a \emph{human} hub user 
triggers irreversible polarization.

Our preliminary work identified a ``backfire effect,'' in which the opinion
shift of moderately opinionated hub users toward extremes, modeled as manipulation to
isolate the structural mechanism, triggers a repulsive backlash across
the network~\cite{Takeda2026}. 
However, these preliminary findings did not systematically
examine the mechanism, boundary conditions, or robustness of this phenomenon---
specifically, how the effect emerges structurally and
how its emergence depends on the
interplay between social modeling parameters and network topology.
The present study addresses this gap by systematically varying agent stubbornness,
opinion prevalence thresholds, and network structure across heavy-tailed,
high-clustering, and random topologies.
\Black
\par


\section{Computational Social Simulation Model}\label{model}
\subsection{Agent-Based Model (ABM)}
We introduce an ABM designed to simulate the opinion
dynamics on social media platforms. Let $G = (\AgentSet, E)$ be a
directed graph representing the social network structure, where
$\AgentSet = \{1, 2, \dots, N\}$ denotes the set of $N$ agents (i.e.,
social media users), and $E \subseteq \AgentSet \times \AgentSet$ denotes the
set of directed edges representing follow relationships. A directed
edge $(i, j) \in E$ indicates that agent $i$ follows agent $j$. We
define the adjacency matrix $A$ such that $A_{ij} = 1$ if $(i,j)\in
E$, and $A_{ij} = 0$ otherwise, where $A_{ij}$ is the $(i,j)$ element
in the $N\times N$ matrix $A$. The dynamics simulations are executed in discrete
time steps, denoted by $t \in \{0, 1, 2, \dots\}$.
\par

\begin{table}[t]
  \tabcolsep=5pt
  \centering
  \caption{Agent parameters and their value ranges}\label{agentParam}
  \begin{tabular}{llll}
    \toprule
    Parameter & Symbol & Range & Initial value \\
    \midrule
    Opinion & $O_i(t)$ & $[-1.0, 1.0]$ & $\mathcal{N}(0.0, 0.6^2)$ \\
    Posting probability & $\post_i(t)$ & $[0.05, 0.5]$ & $0.1$ \\
    Bounded confidence threshold & $\BoundedConfidence_i(t)$ & $[0.2, 1.0]$ & $1.0$\\
    \bottomrule\\
  \end{tabular}
\end{table}

Agent $i \in \AgentSet$ is characterized by the internal parameters listed in Table~\ref{agentParam}. 
Its opinion $O_i(t) \in [-1, 1]$ represents the normalized 
stance of $i$ on a specific issue at time $t$. 
\dgreen{Although opinions are inherently multidimensional, 
many studies in computational social science focus on a 
single salient issue and model opinions along a one-dimensional
spectrum~\cite{Sasahara2021,Uthsav2020}. Under this frameworks, $O_i(t)$
represents agent $i$'s stance toward a specific 
controversial topic, such as support for or 
opposition to a particular \emph{hashtag} on social media. 
Similar issue-specific representations are 
commonly employed in empirical analyses of online 
political communication and polarization~\cite{Conover2011,Toriumi2024}.
}
Boundaries $-1$ and $1$ correspond to the most extreme 
opposing views, and $0$ represents a neutral or moderate 
attitude. We denote the initial opinion $O_i(0)$ as 
the agent's \emph{intrinsic opinion}, 
representing their baseline belief that may 
influence their dynamics over time~\cite{Friedkin1990}, 
and set it to random numbers drawn from a Gaussian
distribution with mean $\mu=0.0$ and standard deviation
$\sigma=0.6$, clipped to the valid range $[-1, 1]$.
To establish a bipolar initial condition, the two highest-degree agents in the initial network are assigned intrinsic opinions of $+1.0$ and $-1.0$, respectively, prior to the simulation.
\par


The posting probability $\post_i(t)$ defines the likelihood
that agent $i$ posts at time $t$. This parameter corresponds to the
concept of \emph{intensity}, which is often used in social simulations.
In our model, posting intensity ($\post_i(t)$) is
defined independently of opinion extremeness ($|O_i(t)|$). That is,
activity is not proportional to radicalization. This independence
allows us to capture heterogeneous behavioral patterns, including
the phenomenon in which moderate users intensify their posting activity. As described in Section \ref{Intro}, \emph{bounded confidence} captures the sociopsychological hypothesis that
individuals are unlikely to engage with opinions that
deviate significantly from their own. Unlike classical bounded confidence
models, in which the threshold directly gates the opinion-averaging
step~\cite{Deffuant2000,Hegselmann2002}, in our model the opinion update
(Eq.~\eqref{FJmodel}) averages over all posts currently in the feed. This
reflects the assumption that a user is inevitably influenced by any information
they actually read: once a post enters and is read from the feed, it affects the
reader's opinion regardless of how far it lies from their own. Bounded confidence
therefore does not filter the opinion update itself; instead it governs the
tie-level behaviors (reposting, following, unfollowing, and confidence
decline) and thereby shapes the opinion dynamics indirectly, by determining
which posts enter the feed in the first place. We further
introduce a \emph{dynamic} bounded confidence parameter $\BoundedConfidence_i(t)$
rather than a static threshold, modeling the adaptive evolution of user
openness to conflicting views through interactions on the platform.
The evolution of these parameters is
governed by the social update mechanisms, as described next.
\par



\subsection{Modeling Social Adaption} 
In general, humans tend to adapt their internal state to their environment. We model these mechanisms based on well-established computational theories~\cite{Friedkin1990,Deffuant2000,Hegselmann2002,Sasahara2021}. Each
time an agent accesses the platform with probability $P_a$, it updates
the internal parameters according to the rules described below. 
Hereafter, we call a post $p_j$ from agent $j$ \emph{comfortable} for agent $i$
when $|O_i(t) - O(p_j)| \leq \epsilon_{\min}$, where $\epsilon_{\min} = 0.2$ is the
minimum bounded confidence threshold (i.e., the lower bound of $\BoundedConfidence_i(t)$);
this notion is used by the social-reinforcement rule below. The dynamic threshold
$\BoundedConfidence_i(t)$ itself instead governs the tie-level behaviors
(reposting, following, unfollowing, and confidence decline), as specified in each rule.
\par



\begin{description}
\item[{\bf Social Influence.}] Agents read all posts in their feeds
  and update their opinions based on posts in feeds and their intrinsic
  opinions. Agent $i$ updates its opinion at time $t$, influenced by
  the opinion of post $p_j$ in its feed $F_i(t)$, as follows:
  %
  \begin{equation}\label{FJmodel}
   O_i(t) = \Stubbornness \cdot O_i(0) + (1 - \Stubbornness) \cdot \frac{\sum_{p_j \in F_i(t)} O(p_j)}{|F_i(t)|},
  \end{equation}
  %
where $\Stubbornness$ denotes the agent’s \emph{stubbornness},
representing the extent to which agents adhere to their intrinsic
opinions $O_i(0)$ rather than being influenced by others~\cite{Friedkin1990}. Note that
posts in an agent's feed originate from those posted or reposted by
users whom the agent follows; thus, the agent may encounter posts that lie
beyond its current bounded confidence unless it unfollows or blocks the creator.

\Dgreen
\item[{\bf Social Reinforcement.}]
 Agents increase or decrease
  their posting probability (or intensity) based on the
  ratio of comfortable posts in their feeds. Formally, let $C_i(t)=
  |K_i(t)|/|F_i(t)|$ be the ratio of comfortable posts in agent $i$'s
  feed at time $t$, where $K_i(t)=\{p_j \in F_i(t) \mid |O_i(t) - O(p_j)|
  \leq \epsilon_{\min}\}$ denotes the set of comfortable posts, as defined above. The posting
  probability is then updated by
  %
  \begin{equation}\label{eq:SocialReinforcement}
    \post_i(t) =
    \begin{cases}
      \min(\post_i(t-1) \cdot (1 + \delta_{\post}),\, 0.5)
        & \text{if } C_i(t) > \ThofOpPreval,\\[2pt]
      \max(\post_i(t-1) \cdot (1 - \delta_{\post}),\, 0.05)
        & \text{else if } C_i(t) \leq 1-\ThofOpPreval,\\[2pt]
      \post_i(t-1) & \text{otherwise,}
    \end{cases}
  \end{equation}
  %
where $\delta_{\post}$ ($> 0$) denotes the parameter controlling the
posting intensity adjustment rate, and $\ThofOpPreval \in (0,1)$ denotes
the {\em threshold for opinion prevalence}. An agent intensifies its posting
when comfortable content locally prevails ($C_i(t) > \ThofOpPreval$) and
reduces it when comfortable content is locally scarce ($C_i(t) \leq 1-\ThofOpPreval$).
The two conditions are evaluated in order, so the increase rule takes precedence;
consequently, when $\ThofOpPreval > 0.5$, an intermediate band
$1-\ThofOpPreval < C_i(t) \leq \ThofOpPreval$ leaves the posting probability unchanged.

\item[{\bf Social Reward.}] In addition to the conformity-based reinforcement above,
  agents receive a posting incentive proportional to the social acceptance of their posts.
  In our implementation, this reward is applied at the \emph{start} of each platform access, 
  reflecting the reposts accumulated since the agent's previous access, prior to reading the current feed.
  When agent $i$'s posts are reposted by others, $i$ updates its posting probability
  by a reward $r_i$ based on the principle of diminishing marginal utility~\cite{Hermann1983}:
  %
  \begin{equation}\label{eq:SocialReward}
    r_i = \alpha \cdot \ln(1 + N^{l}_i),
  \end{equation}
  %
  where $\alpha > 0$ is a tuning parameter and $N^{l}_i$ is the number of reposts received
  since $i$'s previous platform access. The posting probability is then updated as
  $\post_i(t) \leftarrow \min(\post_i(t) \cdot (1 + r_i),\ 0.5)$.
  This logarithmic formulation reflects the empirically observed saturation of social
  validation effects at high repost counts.
\Black

\item[{\bf Confidence Decline.}]
\dgreen{The update of the bounded confidence threshold reflects the asymmetric
  human tendency for negative experiences to carry greater psychological weight
  than positive ones~\cite{Vallone1985,Baumeister2001}; accordingly, exposure to
  conflicting content erodes openness whereas comfortable content does not
  symmetrically restore it, causing $\BoundedConfidence_i$ to evolve monotonically
  downward. Concretely, agents decrease their bounded
  confidence for \emph{each} post they read that lies beyond their current bounded confidence in a given step.
  Formally, $\BoundedConfidence_i$ is decremented multiplicatively by $\delta_{\BoundedConfidence}$ for each post $p_j \in F_i(t)$ that satisfies $|O_i(t) - O(p_j)| > \BoundedConfidence_i$, where $\BoundedConfidence_i$ takes the value already updated by preceding posts within the same step:
  %
  \begin{equation}
    \BoundedConfidence_i \leftarrow \max(\BoundedConfidence_i \cdot \delta_{\BoundedConfidence},\ 0.2)
    \quad \text{for each } p_j \in F_i(t) \text{ s.t. } |O_i(t) - O(p_j)| > \BoundedConfidence_i.
  \end{equation}
  %
  Thus, a single step may produce multiple successive decays if the feed contains multiple such posts.
  The parameters $0 < \delta_{\BoundedConfidence} < 1$ denotes the multiplicative decay rate of confidence in the surrounding environment.
  }
\end{description}



\subsection{Modeling Behaviors on Social Media}
As this study specifically focused on the dynamics within social
media, we introduce the following characteristic behaviors: 
post, repost (like), follow, and unfollow (block).

\begin{description}
\item[{\bf Post.}] Every time agent $j\in\AgentSet$ accesses the
  platform, it creates a post $p_j$ reflecting its current opinion
  with probability $\post_j(t)$. If $j$ creates a post, it is added to
  the feeds of agents who follow $j$; each agent reads at most $\maxfeed$ posts per platform access, in reverse-chronological (newest-first) order, reflecting cognitive limits on information intake.

\item[{\bf Repost (Like).}] When agent $i\in\AgentSet$ views a post in its
  feed, $i$ reposts each post that lies within its bounded confidence
  (i.e., $|O_i - O_{p_j}| <
  \BoundedConfidence_i$) independently with probability $P_r$. Posts are distributed across the social
  media network via this repost function; posts reposted by an agent
  are added to the feeds of its followers. Note that the ``Like''
  function on real-world social media services has a distribution effect
  similar to that described here; therefore, we do not distinguish
  between Like and Repost in the model for simplicity.

\item[{\bf Follow.}] After reading the feed, agent $i$ collects the creators of all posts within its bounded confidence (i.e., $|O_i - O_{p_j}| < \BoundedConfidence_i$) with no existing follow connection as candidates; with probability $P_f$, $i$ randomly selects and follows one candidate.
The maximum number of agents that any agent can follow is fixed at $\maxfollower$ throughout the experiments; in this study, $\maxfollower=10$ based on the concept of limited attention. 
Although users may follow hundreds of other users, the number 
of information sources that can effectively influence them 
at any given time is constrained by human cognitive limits 
and information overload~\cite{Goncalves2011,Weng2012}.
Accordingly, when an agent attempts to follow a new agent while already at $\maxfollower$, it unfollows the single currently-followed agent whose opinion is farthest from its own.

\item[{\bf Unfollow (Block).}] Conversely, if a post lies beyond an agent's bounded confidence (i.e., $|O_i - O_{p_j}| > \BoundedConfidence_i$), the agent unfollows the creator with probability $P_u$, and rejects all future posts. If no connection exists, the agent blocks the creator with probability $P_u$, and rejects all future posts from the creator.
\end{description}
\par

\dgreen{While all agents share the same behavioral rules, heterogeneity 
arises endogenously through individually evolving state variables, 
distinct initial opinions, and the agents' positions; for instance, hub
nodes and peripheral nodes experience qualitatively
different interaction dynamics even under identical rules.}


\subsection{Modeling Manipulation}
\Dgreen
The primary objective of this study is to 
understand the nonlinear dynamics in which an influential \emph{hub} 
user starts to shift their opinion to an extreme end. Accordingly, we manipulated a single hub user to adopt the targeted opinion direction in our simulations. The manipulation target was selected based on ideological position and structural importance in the network. 
After the social media network had evolved sufficiently 
to a steady state, 
agents with neutral opinions (i.e., $|O_i| < 0.2$) and 
follower counts of at least 10\% of the total population (i.e., $\geq 100$ agents for $N = 1000$) 
were identified as candidates. One agent was then randomly 
selected from this pool as the manipulation target. 
If no qualifying candidate was identified, the simulation run was excluded from the results.
\par

We modeled this manipulation as a gradual opinion shift: 
each time the target agent accesses the platform, their opinion is 
incremented by a fixed displacement $\Delta = 0.01$ in the target direction, 
bypassing the social influence update in Eq.~\eqref{FJmodel}, 
until it reaches an extreme value ($+1.0$ or $-1.0$). 
The target's posting probability is simultaneously fixed 
at the maximum value $\rho_{\max} = 0.5$ to reflect 
the heightened activity of a manipulated influencer. 
Under platform access probability of $P_a = 0.1$, the expected number of time steps for the target to 
reach the extreme is approximately 1,000 steps given that the target's opinion is in [-0.2, 0.2]. 
In contrast to the instantaneous shift modeled in our preliminary study~\cite{Takeda2026}, 
this gradual formulation captures 
more realistic process of opinion radicalization. 
Although highly simplified, this single agent setting sufficiently isolates the structural mechanism driving the backfire effect.
This formulation abstracts away the operational mechanism 
by which a hub's opinion shifts, whether through spontaneous radicalization or
coordinated external manipulation, and instead 
isolates the structural consequences of this process.
\Black
\par


\section{Experiments and Results}
\subsection{Experimental Settings} 
\Dgreen
To evaluate the robustness and identify the boundary conditions of the backfire effect, we conducted a sensitivity analysis across five network configurations generated by four distinct models: the \emph{connecting nearest neighbor with random linkage} (CNNR)~\cite{Yuta2007}, the \emph{Holme--Kim} (HK)~\cite{Holme2002}, the \emph{Watts--Strogatz} (WS)~\cite{Watts1998}, and the \emph{Erd\H{o}s--R\'{e}nyi} (ER)~\cite{Erdos1959} models.

The CNNR model extends the original CNN topology~\cite{Vazquez2003} 
by incorporating randomness to simulate human relationship formation 
on social media platforms, reproducing preferential attachment, 
hierarchical clustering, and heavy-tailed degree distributions.
For the HK network, two different conditions were simulated in an analysis, 
namely $p_t = 0$ and $0.3$.
The HK model combines preferential attachment with a triangle-closing 
mechanism controlled by the parameter $p_t$: setting $p_t = 0$ 
recovers pure preferential attachment (yielding a network equivalent 
to the Barab\'{a}si--Albert model~\cite{Barabasi1999} with 
low clustering), while $p_t = 0.3$ additionally closes each newly 
added edge into a triangle with probability $p_t$, generating a 
heavy-tailed network with substantially higher clustering. 
These two HK variants thus allow us to disentangle the contributions of a 
heavy-tailed (hub-dominated) degree structure and local clustering 
within the same generative framework. 
The WS model constructs a small-world network by rewiring a regular
ring lattice with rewiring probability $\beta$, resulting in a topology that exhibits high clustering and short average path lengths but lacks a heavy-tailed degree distribution. The ER model serves as a random baseline lacking both hub structure and clustering. The generation parameters and their resulting average network metrics are summarized in Table~\ref{NWMetrics}.
\par

Crucially, to ensure that any difference in the backfire effect across these
topologies reflects network \emph{structure} rather than a confound of overall
connectivity, we calibrated each model's generation parameter so that all five
networks share a common undirected mean degree $\langle k\rangle = 2E/N \approx 4$
(Table~\ref{NWMetrics}). The topologies therefore differ in their degree
\emph{distribution}, clustering, and path structure while holding edge density
fixed, so that mean degree cannot account for the topology contrast reported
below. As Table~\ref{NWMetrics} shows, under this alignment CNNR and both HK
variants retain heavy-tailed in-degree distributions (power-law exponent
$\approx 2.0$--$2.3$, estimated by maximum likelihood with the \texttt{powerlaw}
package), whereas WS and ER do not. 
\par


The parameters and their values are listed in Table~\ref{ExpParam}.
The intensity adjustment rate $\delta_{\post}$ and confidence decay rate $\delta_{\BoundedConfidence}$ were set based on qualitative reasoning about known social dynamics: bounded confidence erodes gradually under sustained exposure to opposing views, and posting intensity adjusts incrementally in response to social feedback. The empirical adequacy of the resulting model is assessed independently in Section~\ref{sf_validation}.
Each simulation run involved $1{,}000$ agents over $40{,}000$ steps, 
with the manipulation phase introduced at step $20{,}000$. 
The code is publicly available on GitHub.\footnote{\url{https://github.com/tooooomoya/manipulation-backfire-model}}
\par


\begin{table}[t]
  \tabcolsep=5pt
  \centering
  \caption{Experimental parameters}\label{ExpParam}
  \begin{tabular}{lll}
    \toprule
    Parameter & Symbol & Value \\
    \midrule
    Intensity Adjustment Rate & $\delta_\post$ & $0.1$ \\
    Social Reward Tuning Parameter & $\alpha$ & $0.01$ \\
    Confidence Decay Rate & $\delta_\epsilon$ & $0.99$ \\
    Maximum Number of Followers & $\maxfollower$ & $10$\\
    Maximum Feed Size & $\maxfeed$ & $10$\\
    Threshold of Opinion Prevalence & $\ThofOpPreval$ & Variable\\
    Agent's Stubbornness & $\Stubbornness$ & Variable\\
    Probability of Accessing the Platform & $P_a$ & $0.1$\\
    Probability of Reposting & $P_r$ & $0.25$\\
    Probability of Following & $P_f$ & $0.1$\\
    Probability of Unfollowing & $P_u$ & $0.1$\\
    \bottomrule
  \end{tabular}
\end{table}

\begin{table}[t]
  \tabcolsep=5pt
  \centering
  \caption{Network metrics of the initial networks (mean over
  $20$ seeds; path length, diameter, and the power-law exponent over $5$ seeds).}\label{NWMetrics}
  \begin{tabular}{llllll}
    \toprule
    Network topology & CNNR & HK ($p_t{=}0.0$) & HK ($p_t{=}0.3$) & WS & ER \\
    \midrule
    Number of agents               & $1000$  & $1000$  & $1000$  & $1000$  & $1000$ \\
    Average degree  & $4.0$  & $4.0$  & $4.0$  & $4.0$  & $4.0$  \\
    Average path length      & $5.06$ & $3.58$ & $3.26$ & $8.78$ & $5.11$ \\
    Diameter                 & $14$   & $7$    & $7$    & $17$   & $10$   \\
    Power-law exponent       & $2.26$ & $2.13$ & $2.03$ & ---    & ---    \\
    Average clustering coefficient & $0.22$ & $0.05$ & $0.21$ & $0.18$ & $0.002$ \\
    \bottomrule\\
  \end{tabular}
\end{table}

\subsection{Empirical Validation via Stylized Facts}\label{sf_validation}

Agent-based models of social phenomena cannot typically be validated through 
direct predictive testing against held-out data; 
the established alternative is to verify that the model 
reproduces \emph{stylized facts}, which are qualitative regularities robustly
documented in the empirical literature, rather than fitting specific
numerical targets~\cite{Fagiolo2007}. 
We adopt this approach and assess the pre-manipulation equilibrium of 
our model against three stylized facts drawn from empirical studies 
of Twitter/X, using $n = 30$ complete simulation runs under the baseline
parameter settings ($\ThofOpPreval = 0.5$, $\Stubbornness = 0.8$, and the 
Holme--Kim network with $p_t = 0.3$).

\noindent\textbf{SF1: Apparent polarization.}
The variance in observable post content across users consistently
exceeds the variance in the underlying opinion distribution,
because users with more extreme views post disproportionately
more frequently~\cite{Buntain2025}.
We operationalize this phenomenon as the ratio $\sigma_{\mathrm{post}}^{2} / \sigma_{\mathrm{opinion}}^{2}$,
where $\sigma_{\mathrm{opinion}}^{2}$ is the opinion variance weighted by the
number of users at each opinion position (the distribution of opinion holders)
and $\sigma_{\mathrm{post}}^{2}$ is the same variance weighted by the number of
posts contributed at each position (the distribution of observable content).
To anchor this ratio empirically, we recompute it directly from the
publicly released account-level data of~\cite{Buntain2025}
($n = 762$ surveyed Twitter/X users with self-reported ideological leans on a
$[-3, +3]$ scale, contributing $\sim\!1.96$M posts over 2020--2021).

% \footnote{%
% The ratio is invariant to any affine rescaling of the opinion axis, so the
% specific $[-3,+3]$ coding does not affect its value; it relies only on the
% assumption that the seven self-reported categories are approximately
% equally spaced. As a single-study anchor based on a non-representative,
% liberal-leaning sample, this range should be read as an order-of-magnitude
% reference rather than a population parameter.}

In these data, hyper-partisan accounts post roughly two to three times as
often as moderates, inflating the post-weighted opinion variance relative to
the user-weighted variance by a factor of $1.19$ over the full period and
rising to $1.26$ during the high-salience 2020 election window. This later period represents 
precisely the
kind of high-attention regime our manipulation scenario.
This yields an empirical target range of $[1.19, 1.26]$.
Our model achieves a median ratio of $1.232$ (IQR $[1.185, 1.286]$),
with 100\% of runs exceeding the threshold of 1.0.
The median falls within the empirically derived range, confirming that
the model reproduces the stylized fact that publicly observable
opinion is more dispersed, and thus appears more polarized, than
the underlying opinion distribution.

\noindent\textbf{SF2: Cross-partisan repost fraction.}
Repost (retweet) networks are strongly segregated along partisan lines:
cross-partisan reposting is near-zero, far below the rate expected under
random mixing. Conover et al.~\cite{Conover2011} report a retweet-network
modularity of $0.48$ (vs.\ $0.17$ for the mention network) and find that
cross-ideological retweet links occur at only $3$--$5\%$ of their
randomly-expected frequency, concluding that ideologically opposed users
``very rarely share information from across the divide.''
As an upper reference, Wojcieszak et al.~\cite{Wojcieszak2022} find that
users share content from out-group political elites only about
one-fourteenth as often as from in-group elites ($\approx 7\%$ of partisan
shares crossing the divide).
We operationalize this as the fraction of repost edges crossing the opinion
midpoint among partisan agents. Our model yields a median cross-partisan
repost fraction of $2.1\%$ (IQR $[0.6\%, 4.3\%]$), which sits in the
near-zero regime, falling below the $\approx 7\%$ upper reference and aligning 
consistently
with the strong segregation documented above. This supports the validity of
the dynamic bounded confidence and unfollowing mechanisms in generating
realistic levels of within-group information sharing.

\noindent\textbf{SF3: Echo chamber depth.}
Politically engaged users inhabit \emph{partial} echo chambers: a
majority, but not the entirety, of their information diet is
ideologically aligned. Cinelli et al.~\cite{Cinelli2021} show that on
Twitter a user's leaning is strongly correlated with the average leaning of
their network neighborhood, while Bakshy et al.~\cite{Bakshy2015} find that
only about $28.5\%$ of the hard-news content in a user's feed is
cross-cutting, meaning that roughly $70\%$ is ideologically aligned, and the
partisan retweet communities of Conover et al.~\cite{Conover2011} are
$80$--$93\%$ homogeneous. 
We operationalize echo chamber depth as the within-group comfort rate
(the aligned fraction of the feed) for extreme-opinion agents.
Our model yields a median of $64.3\%$ (IQR $[57.7\%, 70.1\%]$), which sits
just below the empirical aligned-share range for engaged users
($\approx\!65$--$76\%$;~\cite{Bakshy2015}). We therefore read SF3 as a
\emph{directional}, not a point, match: the model reproduces the qualitative
signature of a partial echo chamber characterized by a majority-aligned but far-from-total
information diet, while producing a slightly \emph{shallower} chamber than
observed. If anything, this milder homophily makes the model a conservative
substrate for the backfire effect rather than one that inflates it.

The model quantitatively reproduces SF1 and SF2 and matches the qualitative
direction of SF3 (a partial, slightly shallower echo chamber) thereby providing
empirical credibility for the model's structural assumptions in the pre-manipulation
equilibrium.

\subsection{Appearance of the Structural Backfire Effect}\label{backfire_appearance}

\begin{figure}[t]
  \centering
  \includegraphics[width=0.8\linewidth]{figures/backfire_3group_logresp_A2.png}
  \caption{Appearance of the backfire effect in
  a Holme--Kim network
  ($p_t = 0.3$, $\Stubbornness = 0.8$, $\ThofOpPreval = 0.5$). }
  \label{fig:backfire_example_new}
\end{figure}

Consistent with our previous study~\cite{Takeda2026}, we first confirmed the appearance of the backfire effect, here in a Holme--Kim network with $p_t = 0.3$, by setting $\ThofOpPreval=0.5$ and $\Stubbornness=0.8$.
Figure~\ref{fig:backfire_example_new} shows the time-series evolution 
of a per-group \emph{log-response}. For each opinion camp $g$ 
we plot $\log\!\left[\,v_g(t)\,/\,\bar v_g^{\text{pre}}\,\right]$, which is 
the logarithm of its instantaneous posting volume $v_g(t)$ 
normalized by its own mean volume over the pre-manipulation 
window $t \in [19{,}000, 20{,}000]$. This normalization removes 
the large baseline differences in posting volume between camps, 
so that the manipulation-induced change itself becomes the signal. 
The three camps are defined relative to the manipulation direction: 
the \emph{target} camp (the extreme side toward which the hub is 
driven, $|O_i| > 0.2$ on the target sign), the \emph{opposite} camp 
(the opposing extreme), and the \emph{neutral} 
camp ($-0.2 \le O_i \le 0.2$). To ensure the robustness of the observed 
trends, the curves are averaged across $30$ independent 
seeds and pooled over both manipulation directions ($\pm1.0$), 
resulting in $60$ trials, with shaded bands denoting $\pm1$ standard error. 
The manipulation was initiated at the $20{,}000$th step, 
thereafter gradually shifting the targeted hub's opinion 
toward its assigned extreme in the aim of influencing public opinion 
on that side. We summarize each trial by a scalar backfire metric: writing
$\bar{v}^{\text{post}}_{g}$ for the mean posting volume of camp $g$ over the
post-manipulation window $t \in [39{,}000, 40{,}000]$ (and $\bar{v}^{\text{pre}}_{g}$
for the pre-manipulation window defined above), the per-trial backfire effect is

\begin{equation}\label{eq:backfire}
  b = \log\!\left(\frac{\bar{v}^{\text{post}}_{\text{opp}}}{\bar{v}^{\text{pre}}_{\text{opp}}}\right)
    - \log\!\left(\frac{\bar{v}^{\text{post}}_{\text{tgt}}}{\bar{v}^{\text{pre}}_{\text{tgt}}}\right),
\end{equation}

a difference-in-differences estimator on the log scale, where a positive value of
$b$ indicates that the opposing camp grew more active than the targeted camp
relative to the baseline, confirming the occurrence of the backfire effect. By construction, the
gap between the opposite and target curves in
Figure~\ref{fig:backfire_example_new} is the time-resolved metric $b(t)$, which serves as the
running analogue of the scalar $b$.

\par

As can be seen, during the first half of the simulation, prior
to any manipulation, the two conflicting opinionated camps (target and opposite,
which are statistically symmetric at this stage) grow in posting 
activity by a substantially larger factor than the neutral camp, 
as evidenced by their steeper rise toward the baseline. 
This reflects the steady emergence of polarization: as echo chambers 
form on both extremes, agents on the extremes come to post increasingly 
often relative to the moderate majority. We note that, because each 
curve is normalized to its own pre-manipulation window, 
all three necessarily converge at the onset by construction; 
the informative content is therefore the relative steepness 
of the curves before this point, and, most importantly, their
divergence after it.
\par

In the latter half, although intervening on a neutral hub to drive 
it toward the target extreme would intuitively be expected to 
amplify only the target camp, the opposite was observed. 
Immediately after the onset the neutral camp's posting volume 
collapsed sharply dropping to roughly $30\%$ below its baseline on the
log scale, before partially recovering. Furthermore, dropping after a brief
initial transient in which the target camp ticked up, 
the opposite camp's log-response rose above that of the 
target camp and remained elevated long after the manipulation
finished at around the $21{,}000$th step. Equivalently,
the difference $b(t) = \log[v_{\text{opp}}(t)/\bar v_{\text{opp}}^{\text{pre}}] - \log[v_{\text{tgt}}(t)/\bar v_{\text{tgt}}^{\text{pre}}]$ 
became and stayed positive: the opposing camp grew 
more active, relative to its own pre-manipulation baseline, 
than the very camp the manipulation sought to promote. 
Because this pattern is averaged over both manipulation 
directions and $30$ seeds, it is a symmetric 
effect rather than a one-sided artifact of a particular target sign. 
This sustained, counterintuitive outcome, which runs contrary to the
intended purpose of the intervention, is what we
term the \emph{structural backfire effect}.
\par

\subsection{Mechanism of the Structural Backfire Effect}\label{backfire_mechanism}

Regarding the backfire effect confirmed in Section~\ref{backfire_appearance}, 
we hypothesize a four-link
causal chain: an opinion shift by the targeted moderately opinionated hub (1)~drives its opposing
followers to unfollow it, which (2)~strips the network of the bridge that
connected the opposing camp to diverse information, such that (3)~the opposing camp
collapses into a more pure echo chamber, and (4)~this reinforcement raises the
opposing camp's posting activity until it outgrows the targeted camp resulting in the
backfire effect. Because this plausible narrative is not self-evident, this section measures
each link directly on the canonical backfire baseline (Holme--Kim with
$p_t{=}0.3$, $\Stubbornness{=}0.8$, $\ThofOpPreval{=}0.5$; \rev{using $n{=}74$ trials from
$37$ seeds} $\times$ two manipulation directions).
The per-trial outcome is the
log-response ratio $b$ of Eq.~\ref{eq:backfire}; over these trials the backfire effect is
present and robust (\rev{$\bar b = +0.046$, seed-clustered Cohen's $d = +0.97$, one-sided $t$-test
vs.\ $0$: $p < 10^{-6}$}). As shown later
(Table~\ref{tab:doseresponse}), this baseline is a conservative
operating point: \rev{the effect grows further with the manipulated hub's in-degree.}
All quantities below are derived from simulation snapshots.
To enable comparison across manipulation directions, opinion classes are defined
relative to the manipulated side: the camp containing the manipulated hub is
labeled ``targeted,'' whereas the opposite camp is labeled ``opposing.''
This convention allows results from both manipulation directions to be pooled.
\par

\begin{table}[t]
  \tabcolsep=4pt
  \centering
  \caption{Quantitative evidence for the backfire causal chain (HK
  $p_t{=}0.3$, $\Stubbornness{=}0.8$, $\ThofOpPreval{=}0.5$). Values are mean
  $\pm$ standard error across trials, pooled over both manipulation directions.}
  \label{tab:mech_chain}
  \begin{tabular}{llrr}
    \toprule
    \orange{Causal Chain} & \orange{Quantity} & \orange{pre/onset} & \orange{post/end} \\
    \midrule
    \rev{(1)} & \rev{opposing camp unfollows of hub}      & \rev{$7.8 \pm 0.9$}    & \rev{$68.0 \pm 6.2$}  \\
    \rev{(1)} & \rev{Targeted-camp unfollows of hub}      & \rev{$7.8 \pm 0.9$}    & \rev{$13.8 \pm 1.1$}   \\
    \rev{(1)} & \rev{Opposing frac.\ of hub followers}    & \rev{$0.34 \pm 0.01$}  & \rev{$0.00 \pm 0.00$}  \\
    \rev{(1)} & \rev{Targeted frac.\ of hub followers}    & \rev{$0.34 \pm 0.01$}  & \rev{$0.94 \pm 0.01$}  \\
    \midrule
    \rev{(2)} & \rev{Hub betweenness (giant comp.)}     & \rev{$0.094 \pm 0.015$} & \rev{$0.011 \pm 0.002$} \\
    \rev{(2)} & \rev{$\lambda_2$ (giant component)}     & \rev{$0.261 \pm 0.010$} & \rev{$0.236 \pm 0.010$} \\
    \midrule
    \rev{(3)} & \rev{Opposing camp comfort post rate (aggregate)}     & \rev{$0.393 \pm 0.007$} & \rev{$0.412 \pm 0.007$} \\
    \rev{(3)} & \rev{comfort post rate, hub's former followers (opp.)}& \rev{$0.35 \pm 0.02$}   & \rev{$0.48 \pm 0.02$}   \\
    \rev{(3)} & \rev{comfort post rate, non-followers (opp.)}         & \rev{$0.45 \pm 0.01$}   & \rev{$0.45 \pm 0.01$}   \\
    \rev{(3)} & \rev{Network modularity}                 & \rev{$0.45 \pm 0.01$}   & \rev{$0.48 \pm 0.01$}   \\
    \midrule
    \rev{(4)} & \rev{opposing camp post share}           & \rev{$0.429 \pm 0.008$} & \rev{$0.433 \pm 0.007$} \\
    \rev{(4)} & \rev{Targeted-camp post share}           & \rev{$0.429 \pm 0.008$} & \rev{$0.414 \pm 0.007$} \\
    \bottomrule
  \end{tabular}
\end{table}

\subsubsection{Selective unfollowing}
Unfollowing of the manipulation targeted hub is sharply asymmetric and concentrated after the
intervention (Table~\ref{tab:mech_chain}). In the pre-manipulation window the two
camps shed the hub equally (\rev{$7.8$ unfollows each}), whereas after manipulation
\rev{the opposing camp unfollows it nearly five times as often as the targeted camp
($68.0$ vs.\ $13.8$ per trial; paired $t = 10.9$, $p < 10^{-16}$} for the
opposing camp's pre$\rightarrow$post increase). As a result, the hub's follower
set which is balanced at the onset (\rev{$34\%$ opposing, $34\%$ targeted, $32\%$ neutral})
is purged of opposing followers by the end of the run (\rev{$0\%$ opposing, $94\%$
targeted}), thereby converting the manipulated hub from a cross-cutting broadcaster into
a same-side echo node.
\par

\subsubsection{Loss of the bridging role}
Initially, the hub serves as a structural bridge connecting users with diverse political opinions. 
As manipulation progresses, it gradually loses this bridging role due to the loss of cross-cutting 
follower ties. Consistent with this interpretation, the hub's betweenness centrality, 
measured on the undirected giant component, drops by nearly an order of magnitude 
($0.094 \rightarrow 0.011$), providing strong structural evidence that it no longer functions 
as a bridge. At the global level, the algebraic connectivity ($\lambda_2$) of the giant 
component also decreases ($0.261 \rightarrow 0.236$), indicating a modest reduction in 
overall network cohesion. Although echo chambers have already formed to some extent by 
the onset of manipulation, this additional decline in $\lambda_2$ suggests that the loss 
of the hub's cross-cutting ties further fragments the remaining network.
\par

\subsubsection{echo chamber formation}\label{Chain3}
The decisive evidence for the mechanism is its \emph{specificity}. 
At the camp level, the effect is muted as the opposing camp's mean comfort post rate,
whici is the aligned fraction of the feed defined via $\epsilon_{\min}$
 in the Social Reinforcement rule, rises only modestly
 (\rev{$0.393 \rightarrow 0.412$}). On the other hand, disaggregating by prior
attachment to the hub reveals the mechanism. Among opposing agents, those who
\emph{were following the hub} at the onset see their comfort post rate rise from
\rev{$0.35$ to $0.48$} once they lose it, \rev{converging on the level of} opposing agents who were
\emph{not} following the hub, \rev{who remain flat at around $0.45$.
Structurally, the network's modularity rises modestly over the same interval
($0.45 \rightarrow 0.48$), echoing this localized homogenization.} In other words, the hub was precisely
what kept their opposing followers exposed to cross-cutting content so that severing it
collapses them into homogeneous feeds, while agents who never relied on it are
unaffected. This ties the echo chamber to the hub disconnection rather than to
generic repulsion dynamics.
\par

\subsubsection{From echo chamber to posting activity}
Because the backfire effect is defined on posting volume here, the chain must terminate in activity.
The fraction of agents that increase their posting probability after the
intervention is \rev{higher in the opposing and neutral camps ($0.21$ and $0.21$) than
in the targeted camp ($0.19$), and the realized post share of the opposing camp
rises ($0.429 \rightarrow 0.433$) while the targeted camp's share falls
($0.429 \rightarrow 0.414$)}, which is aligned with an asymmetry that $b > 0$ captures.
\par


\subsubsection{Within-model dose-response}
Although the results show the average effect over the trials it still 
leaves open whether the proposed causal chain
operates trial by trial. We therefore ask whether the mechanism is \emph{graded}:
across trials, does a stronger upstream perturbation produce a stronger backfire?
\rev{Indeed the per-trial backfire $b$ increases monotonically with the
structural prominence or the target in-degree at the manipulation onset of the manipulated hub 
as quantified across the heavy-tailed regime in
Table~\ref{tab:doseresponse}.}
This graded response indicates that the loss of the bridging hub
and its follower-specific echo effect are the dominant within-model
drivers of the observed backfire, providing convergent evidence
that the proposed chain is the operative pathway in the present model.
\par

\subsection{Sensitivity Analysis}
To assess the robustness of the backfire effect and delineate its 
boundary conditions across five network configurations spanning four
structural models (CNNR, HK with $p_t{=}0.0$, HK with $p_t{=}0.3$, WS, and ER), 
we conducted a sensitivity analysis by varying two key parameters: 
the threshold of opinion prevalence ($\ThofOpPreval$) and 
social influence stubbornness ($\Stubbornness$). 
The magnitude of the backfire effect was quantified using 
the same per-trial log response ratio $b$ defined in 
Eq.~\ref{eq:backfire} (see Section~\ref{backfire_mechanism}), 
where a positive value of $b$ indicates that the opposing group 
grew more active in posting (relative to its pre-manipulation baseline)
than the targeted group, the signature of the backfire effect.
Each parameter condition was evaluated over up to 100 simulation runs
(50 independent seeds $\times$ two bidirectional manipulation directions).
\rev{Because the two manipulation directions for a given seed are run on the
same network realization, they are not independent observations.
We therefore first average the per-trial backfire $b$ over the two
directions within each seed, obtaining a single value per seed, and
report the effect size as Cohen's $d$ computed from these seed-level
values.}
\par


Statistical significance was assessed via a one-sample $t$-test
of \rev{these per-seed values} against zero (one-sided, alternative: mean $> 0$).
Rejection of the null hypothesis indicates that the 
opposing group systematically outgrew the targeted group 
following manipulation, confirming the backfire effect.
We report significance at the $5\%$ ($p < 0.05$, $*$), 
$1\%$ ($p < 0.01$, $**$), and $0.1\%$ ($p < 0.001$, $*{*}*$) levels.
\par

Figure~\ref{fig:sensitivity_analysis} presents the results as a series of heatmaps, 
one per network configuration, with agent stubbornness $\Stubbornness$ on 
the horizontal axis and the opinion prevalence threshold $\ThofOpPreval$ on the vertical axis. 
Each cell displays the Cohen's $d$ of the backfire effect for that parameter combination; 
cells where the backfire effect is statistically significant are annotated with asterisks. 
This layout allows direct visual comparison of the parameter regions where 
each network is susceptible to the backfire effect.
\par


\begin{figure}[t]
  \centering
  \includegraphics[width=\linewidth]{figures/sensitivity_heatmap.png}
  \caption{Sensitivity analysis of the backfire effect across five network configurations.}
  \label{fig:sensitivity_analysis}
\end{figure}


\Black
Figure~\ref{fig:sensitivity_analysis} shows that the backfire effect 
is a robust phenomenon with respect to $\ThofOpPreval$ across the 
entire tested range ($0.10$ to $0.90$). In the heavy-tailed networks 
(CNNR and HK), once stubbornness was at least
moderately high ($\Stubbornness \geq 0.7$),
Cohen's $d$ remained positive and predominantly significant across the 
full $\ThofOpPreval$ range, indicating that the effect persisted 
regardless of $\ThofOpPreval$. Conversely, the dependency on $\Stubbornness$ is
the dominant pattern: \rev{in the HK networks the backfire effect
became statistically significant already at moderate stubbornness
($\Stubbornness \geq 0.6$), and in CNNR from $\Stubbornness \geq 0.7$,
with effect sizes increasing monotonically up to $\Stubbornness = 0.9$
in all heavy-tailed networks while at the lowest stubbornness, 
($\Stubbornness = 0.5$) Cohen's $d$ was near zero or
negative and never significant.}
Furthermore, in ER and WS networks, which lack heavy-tailed
degree distributions, the backfire effect was
essentially absent. \rev{Cohen's $d$ stayed near zero with no
coherent dependence on either $\Stubbornness$ or $\ThofOpPreval$, and
only isolated cells reached marginal significance ($3$ of $25$ in ER,
$1$ of $25$ in WS, all at the $p<0.05$ level only), at the rate
expected under multiple comparisons. Critically, neither network
exhibits the systematic, stubbornness-gated backfire region that
characterizes the heavy-tailed topologies.}
These results suggest that a heavy-tailed (hub-dominated)
degree distribution is a critical boundary condition for this backfire phenomenon.
\par

\rev{To isolate the manipulated user's own hub status from the heavy-tailed topology
in the aggregate, we measure a graded dose--response on the continuous target
in-degree across the backfire-active regime (the heavy-tailed networks at
$\Stubbornness \geq 0.7$; $45$ conditions, $4{,}294$ trials). Binning by
within-condition in-degree quartile (Table~\ref{tab:doseresponse}), the per-trial
$b$ rises monotonically from $\bar b = +0.018$ (bottom) to $+0.149$ (top), and
correlates positively with in-degree within every condition (Spearman $\rho > 0$ in
$45/45$; pooled within-condition $\rho = +0.27$, $n = 4{,}294$, $p < 10^{-70}$) while ER
and WS, whose most-connected nodes are not genuine hubs, show no effect at all.
Backfire is thus governed by the structural prominence of the manipulated hub,
acting through the channels documented above (Section~\ref{backfire_mechanism}), consistent with the literature
identifying centrality as the primary determinant of systemic
influence~\cite{Kitsak2010,Kempe2003}.}
\par

\begin{table}[t]
  \tabcolsep=6pt
  \centering
  \caption{Dependence of the structural backfire effect on the in-degree of the
  manipulated hub, pooled across the backfire-active regime (heavy-tailed networks
  HK with $p_t{\in}\{0.0,0.3\}$ and CNNR, at $\Stubbornness \geq 0.7$; $45$
  conditions, $4{,}294$ trials).}
  \label{tab:doseresponse}
  \begin{tabular}{lrrrr}
    \toprule
    \rev{Target in-degree} & \rev{$n$} & \rev{mean $b$} & \rev{Cohen's $d$} & \rev{$p$} \\
    \midrule
    \rev{All targets (baseline)} & \rev{$4294$} & \rev{$+0.059$} & \rev{$0.66$}  & \rev{$<10^{-160}$} \\
    \midrule
    \rev{Q1 (lowest)}  & \rev{$1124$} & \rev{$+0.018$} & \rev{$0.42$} & \rev{$<10^{-20}$} \\
    \rev{Q2}           & \rev{$1038$} & \rev{$+0.026$} & \rev{$0.57$} & \rev{$<10^{-30}$} \\
    \rev{Q3}           & \rev{$1096$} & \rev{$+0.046$} & \rev{$0.80$} & \rev{$<10^{-59}$} \\
    \rev{Q4 (highest)} & \rev{$1036$} & \rev{$+0.149$} & \rev{$1.24$} & \rev{$<10^{-100}$} \\
    \bottomrule
  \end{tabular}
\end{table}

In summary, the backfire effect is robust against variations
in the threshold of opinion prevalence, but is sensitive to 
agent stubbornness, being most pronounced in heavy-tailed 
networks (CNNR and HK), in which agents maintain a
high resistance to opinion change. \rev{The within-regime dose--response on target
in-degree further confirms that the structural prominence of the
manipulated user is the proximate driver of the effect.}
\Black
\par


\subsection{Discussion}
Within the scope of our model, these results identify two
boundary conditions for the backfire effect:
(1) the heavy-tailed (hub-dominated) nature of
the network structure and (2) a high degree of agent stubbornness
(or resistance to opinion change).
Both conditions are characteristic of many modern social media
platforms, which suggests, though does not by itself establish, that
such networks may be susceptible to the effect; we return to the
caveats on external validity in the Conclusion.
\par

We posit that the first condition, 
a heavy-tailed (hub-dominated) degree distribution, 
is essential because the backfire mechanism relies 
on the disproportionate influence of hub users. 
In social media ecosystems, ``influencers'' with
enormous followings coexist with the vast majority 
of users who have few connections, forming a 
heavy-tailed degree distribution. This structural
disparity creates fragile connectivity;
when a hub acts as a bridge between opposing views,
their change in state can trigger a cascade of
disconnections, isolating followers into echo chambers, as the causal-chain
analysis in Section~\ref{backfire_mechanism} clarifies.
The consistent susceptibility of both CNNR and HK networks, despite their differing clustering properties,
further confirms that this hub-dominated
degree structure, rather than local clustering, is the decisive structural condition.
\par

Regarding the second condition, we infer that
high stubbornness is a prerequisite for the
formation and maintenance of echo chambers. 
As discussed in Section~\ref{backfire_mechanism}, 
the backfire effect is intrinsically linked 
to the echo chamber formation mechanism. 
When agents hold their opinions firmly, 
they are more likely to reject conflicting information 
from opinion leaders and disconnect from them 
rather than conform to them. Furthermore, our results
indicate that the backfire effect strengthens monotonically
with stubbornness over the tested range: 
in the heavy-tailed networks the mean Cohen's $d$ 
rose steadily from near zero at $\Stubbornness = 0.6$ to 
its maximum at $\Stubbornness = 0.9$, 
with no sign of a turning point within $\Stubbornness \leq 0.9$. 
This implies that a network becomes increasingly vulnerable 
as users hold their intrinsic beliefs more firmly, 
so long as they retain residual susceptibility to 
social influence. Our findings suggest that this 
specific psychological trait, when coupled with 
the network structure, significantly amplifies the 
systemic vulnerability to the backfire effect triggered 
by sudden extreme stance expressions from moderate influencers.
\par

In particular, our results indicate that the threshold 
for opinion prevalence $\ThofOpPreval$, which indirectly 
represents agent conformity to local opinions via 
Eq.~\eqref{eq:SocialReinforcement}, is not the primary 
driver of the backfire effect. Intuitively, one might 
assume that a higher sensitivity to social pressure 
(high $\ThofOpPreval$) would lead to more intense
activation, and thus to a stronger backlash.
However, as discussed in Section~\ref{backfire_mechanism}, 
the backfire effect that we have discussed
in this paper is driven not so much by the magnitude
of user activity but rather by the structural 
disconnection of ties (unfollowing). 
Consequently, the phenomenon remains robust, even when 
agents are not highly sensitive to social pressure.
\par

\Dgreen
The temporal dynamics of the backfire effect 
also merit discussion. As shown in Section~\ref{backfire_appearance}, 
the post-manipulation trajectory exhibits two distinct phases. 
Immediately after the onset, the target camp shows a brief 
transient uptick in posting activity, the expected short-lived
response to the hub's initial opinion shift. However, 
this transient is quickly overtaken: within a short window 
the opposing camp's log-response surpasses that of 
the target camp and \emph{remains} elevated for 
the remainder of the run. Crucially, therefore, 
the backfire effect itself is not transient; it is 
the target-camp spike that is short-lived, and 
what persists is the sustained dominance of the opposing camp. 
This sustained behavioral signal is further compounded 
by the underlying structural damage. The network
fragmentation induced by the backfire effect, most
notably the loss of cross-partisan
ties between the hub and their former opposing
followers, persists for the remainder of the simulated horizon,
well after posting volumes stabilize, and these ties do not reform
within our observation window. This indicates that the structural consequences
are at least as durable as the already-persistent surface-level behavioral signal.
\par


\section{Conclusion}
\Dgreen
The structural backfire effect is the counterintuitive finding that, in online
social networks, 
opinion shifts by a moderate influential hub toward an extreme
stance energize the opposing camp rather than the targeted one. This paper
establishes the effect as a robust phenomenon and identifies the mechanism that
produces it. Building on a
preliminary observation, and using an ABM grounded in well-established
social-interaction protocols with the manipulation modeled as an incremental opinion shift,
we provide
direct mechanistic evidence for the causal chain underlying the effect, delineate
its boundary conditions through a sensitivity analysis spanning four
network-generating models and a two-parameter grid of agent stubbornness and
opinion prevalence threshold, and establish the model's behavioral credibility
against empirical stylized facts drawn from real-world datasets.
\par


The backfire effect proved robust to the threshold of opinion prevalence (the
local conforming rate) across its entire tested range, indicating that the
phenomenon is driven by structural reorganization rather than by local conformity
pressure. Its emergence instead hinges on two boundary conditions: a heavy-tailed,
hub-dominated degree distribution and a relatively high degree of agent adherence to
intrinsic opinions. The effect strengthened monotonically with stubbornness and
was essentially absent in networks lacking a heavy tail, which is confirmed by
restricting the analysis to genuinely high-degree targets
sharpened this pattern, expliciting that the structural prominence of the shifted
hub is its proximate driver. Direct measurement of the causal chain further showed
that the effect operates through echo chamber deepening via hub loss: the
manipulated hub is selectively abandoned by their opposing followers, severing the
cross-partisan bridge it provided and collapsing those former followers, but not
agents who never relied on it, into homogeneous feeds, thereby affecting their echo chamber formation.
\par


Several limitations bound the scope of these conclusions. First, our claims are
internal to the model: the backfire effect is established as a robust property of
a particular class of opinion dynamics models rather than as a directly observed
empirical regularity. Although we anchored the pre-manipulation equilibrium to
three stylized facts, this evidence is qualitative and partly directional, 
is drawn from a modest number of real dataset studies
covering specific platforms and time windows, and does not validate the
post-manipulation dynamics themselves, for which comparable ground-truth data
remain scarce. This asymmetry is itself informative: whereas the pre-manipulation
equilibrium can be anchored to established stylized facts, to the best of our knowledge, 
comparable empirical regularities have not yet sufficiently been documented for the post-backfire regime, the very
window in which the structural backfire effect unfolds. We therefore treat this gap
not merely as a caveat but as a methodological opening for future work: it
identifies a concrete need to establish post-intervention stylized facts and to
specify observable structural signatures of the mechanism we describe (for example,
hub-loss-driven fragmentation or echo chamber-deepening indices) that could be
measured in longitudinal platform traces, thereby offering a stepping stone toward
the empirical validation of post-manipulation dynamics. Second, the model is
simplified for tractability. 
For instance, opinion is abstracted as a single one-dimensional issue and
the feed is governed solely by follow ties; we deliberately omit algorithmic
recommendation and curation, which demonstrably shape exposure on contemporary
platforms and could either amplify or dampen the structural cascade we describe.
Third, the manipulation is reduced to a single, exogenously driven hub whose
opinion is displaced deterministically. This design isolates the structural
mechanism but does not capture coordinated multi-actor influence operations or the
endogenous heterogeneity of real radicalization trajectories. 
\par


Addressing these limitations constitutes a natural and important direction for
future work, in particular through comparative validation against empirical traces
of real-world influence operations, the incorporation of algorithmic curation, and
richer multi-actor manipulation scenarios. Notwithstanding these caveats, our
findings consistently identify the joint role of hub-dominated structure and user
stubbornness in rendering networks vulnerable to backfire, offering a structural
framework for anticipating when well-intentioned or adversarial interventions on
influential users may perversely deepen polarization, and thereby providing
organizing principles for the design of more resilient web-communication systems.
\par


%%===========================================================================================%%
%% If you are submitting to one of the Nature Portfolio journals, using the eJP submission   %%
%% system, please include the references within the manuscript file itself. You may do this  %%
%% by copying the reference list from your .bbl file, paste it into the main manuscript .tex %%
%% file, and delete the associated \verb+\bibliography+ commands.                            %%
%%===========================================================================================%%


%Added by Sugawara, eliminate the next one line when submission.
%\clearpage
\Black

\backmatter
%\Blue

\iffalse
\bmhead{Supplementary information}
If your article has accompanying supplementary file/s please state so here. 
Authors reporting data from electrophoretic gels and blots should supply the full unprocessed scans for key as part of their Supplementary information. This may be requested by the editorial team/s if it is missing.

\bmhead{Acknowledgements}
Acknowledgements are not compulsory. Where included they should be brief. Grant or contribution numbers may be acknowledged.

\fi

\section*{Declarations}

\subsection*{Funding} This work is partly supported by JST ERATO (Grant Number JPMJER2502), 
by JSPS KAKENHI (Grant Numbers 25K03188 and 22H03906), 
and by the Open Access publication support system of Waseda University.
\subsection*{Conflict of interest}
The authors declare no conflict of interest.
\subsection*{Data availability}
No datasets were generated or analyzed during the current study.
\subsection*{Code availability}
The modified implementations of all algorithms
used in the experiments of this study can be found at https://github.com/tooooomoya/manipulation-backfire-model
\subsection*{Author contribution}
Conceptualization: TT, FT and TS; Methodology: TT;
Implementation and experimental analysis: TT; Writing - original
draft preparation: TT; Writing - review and editing: TS and FT;
Supervision: TS
\subsection*{Corresponding author}
Correspondence to Tomoya Takeda.




\iffalse
  Some journals require declarations to be submitted in a standardised format. Please check the Instructions for Authors of the journal to which you are submitting to see if you need to complete this section. If yes, your manuscript must contain the following sections under the heading `Declarations':

\begin{itemize}
\item Funding
\item Conflict of interest/Competing interests (check journal-specific guidelines for which heading to use)
\item Ethics approval and consent to participate
\item Consent for publication
\item Data availability 
\item Materials availability
\item Code availability 
\item Author contribution
\end{itemize}

\noindent
If any of the sections are not relevant to your manuscript, please include the heading and write `Not applicable' for that section. 


%%===================================================%%
%% For presentation purpose, we have included        %%
%% \bigskip command. Please ignore this.             %%
%%===================================================%%
\bigskip
\begin{flushleft}%
Editorial Policies for:

\bigskip\noindent
Springer journals and proceedings: \url{https://www.springer.com/gp/editorial-policies}

\bigskip\noindent
Nature Portfolio journals: \url{https://www.nature.com/nature-research/editorial-policies}

\bigskip\noindent
\textit{Scientific Reports}: \url{https://www.nature.com/srep/journal-policies/editorial-policies}

\bigskip\noindent
BMC journals: \url{https://www.biomedcentral.com/getpublished/editorial-policies}
\end{flushleft}

\fi

%\bibliographystyle{bst/sn-mathphys-num}
\bibliography{refs}% common bib file

%% if required, the content of .bbl file can be included here once bbl is generated
%%\input sn-article.bbl

\end{document}





